\begin{frame}[t]{Phase 4: A Brief Look at Trees}
\vspace{.42cm}
\lead{Our boosting models learned a sequence of simple corrections.}
\vspace{.28cm}
\begin{definitionbox}
\centering\vspace{.26cm}
{\fontsize{21}{26}\selectfont How does a decision tree\\build one prediction?\par}
\vspace{.26cm}
\end{definitionbox}
\vspace{.48cm}
\begin{takeaway}
\textbf{Next:} Partition the feature space, fit a value in each leaf, and examine what changes when the data change.
\end{takeaway}
\note{Introduce just enough tree mechanics for the coming bagging discussion: a recursive partition, a leaf prediction, a greedy split, and sensitivity to the training data. Earlier boosting corrections included threshold rules and piecewise-constant functions, while the electricity application used spline learners. Trees are a common choice of base learner, not a requirement of gradient boosting.}
\end{frame}
